\section{Decisiones y consideraciones de diseño}

\begin{description}[leftmargin=1.5em,style=nextline]
  \item[Inventario por sucursal.]
  \texttt{cantidadDisponible} pertenece a \texttt{Disponer} y no a
  \texttt{Premio}. Así, un mismo premio conserva una descripción única, pero
  puede tener existencias distintas en cada sucursal.

  \item[Operaciones con identidad propia.]
  Partidas, visitas, recargas, canjes y mantenimientos se modelan como
  entidades porque requieren identificador, atributos temporales y relaciones
  con varios participantes. Esta decisión preserva el historial de operación.

  \item[Separación entre juego, máquina y partida.]
  \texttt{Juego} describe la oferta; \texttt{Máquina}, el equipo físico de
  una sucursal; y \texttt{Partida}, cada ejecución pagada por una tarjeta. La
  separación evita repetir costo y nombre del juego en cada máquina o partida.

  \item[Reglas de puntuación configurables.]
  \texttt{Tipo de juego} se modela como catálogo relacionado mediante
  \texttt{Clasificar}. Sus atributos permiten describir las modalidades y sus
  condiciones sin alterar la estructura del modelo.

  \item[Datos derivados.]
  \texttt{edad}, \texttt{categoríaVIP} y \texttt{bonificaciónVIP} se marcan
  como derivados para evitar valores incompatibles con la fecha de nacimiento,
  el historial de visitas o los puntos obtenidos.

  \item[Tarjeta como medio operativo.]
  La tarjeta concentra saldo y puntos, registra partidas, recibe recargas y
  origina visitas. La relación 1:1 con el cliente expresa la regla de una sola
  tarjeta activa, mientras \texttt{estado} permite conservar antecedentes de
  activación.

  \item[Responsabilidad del personal.]
  La especialización total y disjunta de \texttt{Empleado} permite vincular
  cada operación con el puesto autorizado sin duplicar los datos comunes del
  personal.

  \item[Identidad del empleado.]
  El diagrama subraya \texttt{CURP} y \texttt{RFC}. Se interpretan como claves
  candidatas únicas y heredadas por las subentidades; la implementación
  relacional posterior podrá elegir una como clave primaria y mantener la otra
  con restricción de unicidad.
\end{description}
